\chapter*{摘要}
\addcontentsline{toc}{chapter}{摘要}
医学图像检测与分割是医学人工智能中的基础环节，也是数字病理定量分析、疗效评估和临床辅助决策的重要前提。与自然图像相比，医学图像具有目标尺度小、实例密集、类间差异细微、成像与染色条件变化显著等特点；高质量框标注、实例掩码及全切片区域标注又高度依赖专业医师，获得成本高且难以在多中心、大规模数据上持续扩展。与此同时，传统封闭集模型通常要求训练类别与测试类别保持一致，一旦面对新的器官、细胞类型、疾病亚型或中心分布，往往需要重新标注和训练。如何充分利用大规模视觉语言预训练形成的开放语义知识，并通过提示学习把自然图像先验稳健迁移到医学场景，是降低标注成本、提高跨域泛化能力的重要研究问题。

围绕“低标注成本条件下的医学图像检测、分割与临床量化”这一主线，本文系统研究了自动属性提示、反馈式自训练、视觉文本化图像提示以及肺癌数字病理分析四个相互递进的问题。总体思路是：首先利用语言描述建立未标注医学目标与视觉语言模型之间的语义联系；随后利用模型预测、实例特征和伪标签反馈修正初始检测；进一步将少量支持图像转换为可参与跨模态推理的视觉 token，以补充文本难以表达的形态和染色信息；最后将检测、分割结果汇总为具有病理学意义的区域、细胞和空间统计指标，验证方法对真实临床研究的支撑能力。

首先，针对细胞核类别名称信息量有限、人工提示依赖经验且目标域缺少检测标注的问题，本文提出基于自动属性提示的零样本细胞核检测方法。该方法以视觉问答和属性提取获得颜色、形状等细粒度描述，将其组织为适合视觉语言模型的检测提示，并利用原始预测生成伪标签开展自训练。与仅输入简单类别名的零样本推理相比，属性提示能够提供更明确的视觉定位线索；自训练则把一次性预测转换为可迭代的数据闭环，使模型逐步适应病理图像中的染色变化、密集排列和小目标尺度。该研究证明，在不依赖目标域人工框标注的条件下，视觉语言预训练模型可以被迁移到细胞核检测任务，并为后续无监督优化建立了基础。

其次，针对自动提示中候选属性冗余、提示顺序敏感以及密集细胞核漏检等问题，本文进一步提出面向密集细胞核检测的属性提示增强与自训练知识蒸馏框架。该框架通过属性扩展、相关性排序和模板化组合提高提示的覆盖范围与判别性，并以教师模型输出为伪监督信号训练 GLIP 学生模型，在保持预训练跨模态知识的同时增强目标域定位能力。围绕反馈式优化，本文还研究了高置信度实例框特征作为模板 token 的迭代注入、实例特征与教师—学生自训练的协同作用、中文视觉语言权重与中文医学提示的适配，以及分割头对检测结果向实例掩码扩展的影响。这些补充研究从语义输入、视觉反馈和任务输出三个层面解释了方法发挥作用的条件，也明确了伪标签噪声累积、类别不均衡和提示翻译误差等边界。

再次，针对纯文本提示难以完整表达未见目标外观、支持图像信息又难以直接进入语言条件检测器的问题，本文提出视觉文本化图像提示方法。该方法将少量支持图像中的多尺度视觉特征映射到视觉语言模型已经对齐的文本特征空间，使图像示例以视觉 token 的形式与类别文本共同参与检测。通过多尺度文本化和多阶段融合，模型在不破坏预训练主体结构的前提下同时保留类别语义锚点和支持实例外观信息。本文在自然图像少样本检测、开放集跨数据集迁移以及 MoNuSeg、CoNSeP、CCRCC 等医学病理数据上进行验证，并从定量结果、注意力响应、定性可视化和特征分布等角度分析其有效性。结果表明，视觉文本化能够缓解仅依赖类别名时的语义粒度不足，在跨域、小目标和外观变化明显的场景中形成更稳定的目标响应。

最后，本文将上述低标注成本识别能力应用于新辅助免疫联合化疗后肺鳞癌的数字病理分析，构建由多尺度全切片区域分割、多类肿瘤周围细胞检测、残余活性肿瘤比例计算和空间免疫统计组成的分析流程。该流程在切片和患者两个层面比较 AI 量化与病理医师视觉评估的一致性，并进一步分析残余肿瘤负荷、淋巴细胞、中性粒细胞、内皮细胞等空间特征与无病生存期和总生存期之间的关系。研究结果显示，自动量化能够在跨中心数据中保持较稳定的一致性；相对于单一阈值判断，连续残余肿瘤指标保留了更充分的病理反应信息；将细胞类别与其相对肿瘤边界的距离联合建模，还可以揭示局部炎症生态和广域抗肿瘤免疫之间的差异。

综上，本文形成了一条由“语义提示构造—预测反馈优化—图像提示注入—临床指标量化”组成的技术路线。主要贡献包括：建立无目标域框标注的细胞核零样本检测入口；提出面向密集病理实例的自动属性提示与知识蒸馏机制；实现与视觉语言预训练空间兼容的图像提示表示；并将检测、分割结果转化为可复核的病理量化和预后分析依据。相关研究为视觉语言预训练模型在医学场景中的低成本迁移提供了方法基础，也为开放类别医学图像理解和可解释数字病理研究提供了可扩展的实现路径。

\noindent\textbf{关键词：} 医学图像分析；细胞核检测；实例分割；视觉语言预训练；提示学习；数字病理

\chapter*{Abstract}
\addcontentsline{toc}{chapter}{Abstract}
Medical image detection and segmentation are fundamental components of medical artificial intelligence and are essential for quantitative digital pathology, treatment-response assessment, and clinical decision support. Compared with natural images, medical images often contain small and densely distributed targets, subtle inter-class differences, and substantial variations caused by scanners, institutions, staining protocols, and disease subtypes. High-quality bounding boxes, instance masks, and whole-slide region annotations depend heavily on experienced pathologists and are difficult to scale across large multi-center cohorts. Conventional closed-set models also assume that training and testing categories are fixed. Once a new organ, cell type, disease subtype, or acquisition domain is introduced, extensive re-annotation and retraining are commonly required. Therefore, transferring open semantic knowledge from large-scale vision-language pre-training to medical images through prompt learning is important for reducing annotation cost and improving cross-domain generalization.

This thesis studies four progressively connected problems under the theme of low-annotation-cost medical image detection, segmentation, and clinically interpretable quantification: automatic attribute prompting, feedback-driven self-training, visual textualization for image prompting, and digital pathology analysis of lung squamous cell carcinoma. The general strategy is to first establish semantic links between unlabeled medical targets and a vision-language model through descriptive prompts; then use model predictions, instance features, and pseudo labels to refine initial detection; further convert a small number of support images into visual tokens that can participate in cross-modal inference; and finally aggregate detection and segmentation outputs into region-level, cellular, and spatial statistics with pathological and clinical meanings.

First, a zero-shot nuclei detection method based on automatic attribute prompting is developed to address the limited information carried by class names and the lack of target-domain bounding-box annotations. Color, shape, and other fine-grained attributes are obtained through visual question answering and attribute extraction, and are organized into detection prompts suitable for a vision-language model. Raw predictions are subsequently converted into pseudo labels for self-training. Attribute-aware prompts provide more explicit visual localization cues than simple category names, while self-training turns one-pass inference into an iterative data loop that adapts the model to staining variation, dense nuclei, and small object scales. This study demonstrates that vision-language pre-trained models can be transferred to nuclei detection without manually annotated target-domain boxes.

Second, an attribute-prompt enhancement and self-trained knowledge-distillation framework is proposed for dense nuclei detection. Attribute expansion, relevance-based ranking, and template composition improve the coverage and discriminability of prompt candidates. Predictions from a teacher model are used as pseudo supervision for a GLIP student, enhancing target-domain localization while preserving pre-trained cross-modal knowledge. Additional studies investigate iterative injection of high-confidence box features as template tokens, cooperation between instance-feature feedback and teacher-student training, adaptation of Chinese vision-language weights with translated Chinese medical prompts, and extension from boxes to masks through a segmentation head. These studies clarify why feedback-driven prompting can work and also reveal practical limitations caused by pseudo-label noise, class imbalance, and prompt translation errors.

Third, a visual textualization framework is introduced because purely textual prompts cannot fully express the appearance of unseen targets, whereas image exemplars cannot be directly consumed by a language-conditioned detector. Multi-scale support-image features are mapped into the text feature space already aligned by the vision-language model. The resulting visual tokens are fused with category text during detection, preserving the semantic anchor of the class name while supplementing it with target-domain appearance. Experiments on natural-image few-shot detection, open-set cross-dataset transfer, and medical pathology datasets including MoNuSeg, CoNSeP, and CCRCC are analyzed through quantitative metrics, attention responses, qualitative visualization, and feature distributions. The results show that visual textualization alleviates insufficient semantic granularity and produces more stable target responses under domain shift, small-object conditions, and substantial appearance variation.

Finally, the developed recognition capability is applied to digital pathology of lung squamous cell carcinoma after neoadjuvant immunochemotherapy. A multi-stage pipeline is constructed for multi-scale whole-slide region segmentation, multi-class peritumoral cell detection, residual viable tumor quantification, and spatial immune analysis. AI-derived measurements are compared with visual assessment by pathologists at both slide and patient levels. Residual tumor burden and spatial distributions of lymphocytes, neutrophils, endothelial cells, and other cell types are further analyzed in relation to disease-free and overall survival. The results indicate stable cross-center agreement, demonstrate that continuous residual-tumor measurements preserve more response information than a single threshold, and show that jointly modeling cell identity and distance to the tumor boundary can reveal complementary local inflammatory and broad anti-tumor immune patterns.

In summary, this thesis establishes a coherent technical route from semantic prompt construction, prediction-feedback optimization, and image-prompt injection to clinically interpretable quantification. The contributions provide a practical foundation for low-cost adaptation of vision-language models to medical images, extend open-vocabulary perception to dense pathological targets, and connect detection and segmentation outputs with reproducible pathological measurements and prognosis analysis.

\noindent\textbf{Key Words:} medical image analysis; nuclei detection; instance segmentation; vision-language pre-training; prompt learning; digital pathology
